\scanpage{23}\mspage{II}
$\in S(\Gamma)$ appartient à $\Gamma^{i}_{xy}$ alors $s\in S(\Gamma^{i}_{xy}),\ i=1,2$.
\margin{\tikz[x=.06cm,y=.06cm]{\draw (1,0)--(0,2)--(.2,3)--(-.5,4)--(0,5)--(-.4,6)--(.5,7)--(-.3,7.8)--(-.8,7.2);}}

Si $\Lambda^{1}_{xy}$ et $\Lambda^{2}_{xy}$ sont deux arcs de polygone vérifiant \margin{$\bullet$}$\Lambda^{1}_{xy}\cap\Lambda^{2}_{xy}=\{x,y\}$ on peut obtenir par \unsure{récollement|recollement} selon $\{x,y\}$ un polygone $\Lambda$ de façon évidente.

\margin{\unsure{fin|bien|lin}}%
\uline{Definition 1}. On dit \unsure{qu'un|que'un} polygone $\Gamma$ est un $J$-polygone, ou bien que $\Gamma$ est de Jordan si $\Gamma$ vérifie :

a) $\mathbb{R}^{2}-\Gamma=U\cup V$ avec $U$ et $V$ connexes, $U$ relativement compact et \uline{$\overline{V}$ non-compact}, On pose $U=\mathrm{int}\,\Gamma$, $V=\mathrm{ext}\,\Gamma_{\unsure{\bullet|,|x}}$ ;
\margin{$V$\,! non rel. compact}

b) Si $D_{x}$ est un disque de centre $x\in\Gamma$ qui ne coupe que les arêtes de $\Gamma$ auxquelles $x$ appartient (deux au plus), alors $D_{x}-D_{x}\cap\Gamma=U'\cup V'$, avec $U'\subset\mathrm{int}\,\Gamma$ et $V'\subset\mathrm{ext}\,\Gamma$ connexes\inserted{(i.e. les 2 comp. conn. de $D_{x}-D_{x}\cap\Gamma$ \deleted{\unsure{appartiennent|appartient}} \inserted{\unsure{sont contenues|sont contenus}}, \unsure{l'un|l'une} dans $\mathrm{int}\,\Gamma$), \unsure{l'aut|l'autre} \unsure{ds|de|dans} $\mathrm{Ext}\,\Gamma$\unsure{)|,}}.
\margin{non \uuline{\unsure{vides|vide}}\,!}

\margin{style}%
\uline{Remarque} 2 On voit qu'il suffit qu'un seul disque \deleted{tel} qui ne coupe $\Gamma$ \deleted{ailleurs} que dans les \deleted{\illegible{1}} arêtes auxquelles $x$ appartient vérifie la propriété de b), pour que tout autre \inserted{ayant cette \deleted{\unsure{mê|m}} même propriété} disque \inserted{la} vérifie \inserted{(car \inserted{si $D_{x}\subset D'_{x}$,} $\pi_{0}(D_{x}-D_{x}\cap\Gamma)$ $\longrightarrow$ $\pi_{0}(D'_{x}-\unsure{D'_{x}|D'_{\bullet}}\cap\Gamma)$ est bijectif)}

On pose $R_{\Gamma}=\mathrm{int}\,\Gamma\cup\Gamma$.

\margin{$\cdot$}\uline{Definition} 2. On dit que $s\in S(\Gamma)$ est un \uline{bon sommet} de $\Gamma$ s'il \margin{,}n'est pas colineaire avec les deux sommets $s',s''$ qui sont adjacents à $s$, et si le triangle de sommets $s,s',s''$ ne contient aucun sommet de $\Gamma$ à l'interieur. \inserted{(On \uline{suppose connue} la \uline{notion} de interieur ou de exterieur d'un triangle\unsure{.)|).}}

\margin{\unsure{vaguë|vague}}%
\inserted{\uline{Déf 3}?}

\margin{\tikz[x=.1cm,y=.1cm]{\draw (0,0)..controls (.2,1.2) and (1.3,.4)..(1.5,1.4);}}%
\margin{\rule{.35pt}{1em}}%
Si $E$ est une partie de $\mathbb{R}^{2}$ \inserted{et $s,s'\in\mathbb{R}^{2}$,} on dit que $s'$ est $E$-visible de $s$ si le segment \margin{\tikz[x=.1cm,y=.1cm]{\draw (0,0)--(0,1) (0.4,.65)--(1.5,.65)--(1.5,0);}}ouvert $]s,s'[$ est contenu dans $E$ \deleted{$/s,s'\in\mathbb{R}^{2}/$}\unsure{,|.}

\margin{\rule{.35pt}{1em}}Si $\Gamma$ est un $J$-polygone on dit que $s'$ est visible de $s$ de l'exterieur \inserted{(resp. interieur)} de $\Gamma$ si $s'$ est $\mathrm{ext}\,\Gamma$ (res. $\mathrm{int}\,\Gamma$) - visible de $s$.\margin{\rule{.35pt}{1em}}
% archive sheet number: 22
